El experimento comenzó con la calibración del campo magnético $\mathbf{B}.$ Para esto se hizo uso de 5 bobinas, $A_1,\, A_2,\, A_3,\, B_1,\, B_2$,
, cada una con un reóstato en serie. Dichas bobinas se encargan de generar el campo magnético $\mathbf{B}$ dentro del cilindro donde irá el tubo de rayos catódicos. Cabe recalcar que dicho cilindro se encuentra inclinado de forma tal
que su eje es paralelo a la dirección del campo magnético terrestre, de modo que el campo generado por las bobinas se suma directamente con el terrestre.

Se hizo uso de un multímetro Fluke179 True RMS para medir la corriente total $I$ del circuito y una fuente GW Instek PSH-3020A
para suministrar corriente al circuito, como se observa en la Figura \ref{fig:arreglo}. Para la medición del campo magnético se hizo uso de un gaussímetro LakeShore 425 con una sonda Hall axial.

\begin{figure}[h!]
    \centering
    \begin{tikzpicture}[scale=.9]
        % Paths, nodes and wires:
        \draw (5, 5) to[american inductor, l={$A_1$}] (5, 3);
        \draw (6, 5) to[american inductor, l={$A_2$}] (6, 3);
        \draw (7, 5) to[american inductor, l={$A_3$}] (7, 3);
        \draw (8, 5) to[american inductor, l={$B_1$}] (8, 3);
        \draw (9, 5) to[american inductor, l={$B_2$}] (9, 3);
        \draw (5, 3) to[variable american resistor] (5, 1);
        \draw (6, 3) to[variable american resistor] (6, 1);
        \draw (7, 3) to[variable american resistor] (7, 1);
        \draw (8, 3) to[variable american resistor] (8, 1);
        \draw (9, 3) to[variable american resistor] (9, 1);
        \node[shape=rectangle, minimum width=0.465cm, minimum height=0.465cm](N1) at (6.5, 3){} node[anchor=center] at (N1.text){$R_{A2}$};
        \node[shape=rectangle, minimum width=0.465cm, minimum height=0.465cm](N2) at (7.5, 3){} node[anchor=center] at (N2.text){$R_{A3}$};
        \node[shape=rectangle, minimum width=0.465cm, minimum height=0.465cm](N3) at (5.5, 3){} node[anchor=center] at (N3.text){$R_{A1}$};
        \node[shape=rectangle, minimum width=0.465cm, minimum height=0.465cm](N4) at (8.5, 3){} node[anchor=center] at (N4.text){$R_{B1}$};
        \node[shape=rectangle, minimum width=0.465cm, minimum height=0.465cm](N5) at (9.5, 3){} node[anchor=center] at (N5.text){$R_{B2}$};
        \draw (2, 1) to[ammeter] (4, 1);
        \draw (4, 1) -- (9, 1);
        \draw (2, 2.25) to[american current source] (2, 3.75);
        \draw (2, 1) -- (2, 2.25);
        \draw (2, 3.75) -| (2, 5) -- (5, 5);
        \draw (5, 5) -- (9, 5);
        \node[shape=rectangle, minimum width=1.965cm, minimum height=1.465cm] at (3, 4.25){} node[anchor=center, align=center, text width=2cm, inner sep=6pt] at (3.35, 3.25){GW Instek PSH-3020A};
    \end{tikzpicture}
    \caption{Arreglo experimental para la calibración del campo magnético $\mathbf{B}$.}
    \label{fig:arreglo}
\end{figure}

La calibración consistió en medir el campo $B$ en distintas posiciones dentro del tubo,
 con la ayuda de una vara de madera con el sensor Hall ubicado sobre su punta. 
 Dicha vara era introducida dentro del tubo a distintas profundidades y en cada posición
 se registraban los valores de $B$ e $I,$ luego se variaban los valores de los reóstatos a
 modo de minimizar la variación de $B$ a lo largo del tubo. Una vez obtenido un campo uniforme,
  se procedió a medir el valor del campo $B$ en distintas posiciones del cilindro para
  distintos valores de corriente $I,$ a modo de establecer una relación entra la corriente total $I$
  del circuito y el campo $B.$ Para ver los datalles del proceso ver el Anexo A.

Una vez terminada dicha calibración se procedió a ubicar el tubo de rayos catódicos en la zona de menor variación de campo $B$